\begin{ledgroupsized}[r]{120mm}
\footnotesize
\pstart
\noindent\textbf{Überlieferung:}
\pend
\end{ledgroupsized}
\begin{ledgroupsized}[r]{114mm}
\footnotesize
\pstart \parindent -6mm
\makebox[6mm][l]{\textit{L}}%
Auszüge mit Bemerkungen aus 
\indexN{\textls{Descartes} (Cartesius, des Cartes), René 1596\textendash1650}\textsc{R.~Descartes},
\cite{00035}\textit{Principia Philosophiae}, Pars~II, §§~46\textendash52, Amsterdam 1644, S.\;59\textendash62:
LH~XXXV~10, 12~Bl.~38.
Ein Blatt (ca.~20~x~28~cm.);
\href{\Scan{LH_35_10_12_038wz}}%
{Wasserzeichen}
in der Blattmitte (Papier aus Osterode im Harz); %
alle Ränder beschnitten.
Zwei Seiten;
Bl.~38~r\textsuperscript{o} einspaltig beschrieben.
\pend
\end{ledgroupsized}
\vspace{5mm}
\begin{ledgroup}
\footnotesize
\pstart
\noindent%
\textbf{Datierungsgründe:}	
Die Datierung %
der vorliegenden kommentierten Auszüge aus dem zweiten Teil von
\indexN{\textls{Descartes} (Cartesius, des Cartes), René 1596\textendash1650}Descartes' \cite{00035}\textit{Principia philosophiae}
stützt sich einerseits
auf das für den Zeitraum 1680\textendash1685 mehrfach belegte Wasserzeichen (Papier aus dem Harz),
andererseits
auf die inhaltliche Verwandtschaft mit den Randbemerkungen in Leibnizens Handexemplar der \cite{00035}\textit{Principia}, deren Entstehung ab 1685 datierbar ist, 
und mit Malebranches \cite{02080}\textit{De la Recherche de la vérité}.
Mit den um 1689/1690 während des Italienaufenthaltes angefertigten Exzerpten derselben Descartes-Stelle
(\cite{02091}\textit{LSB} VIII,~3 N.\;74)
ist N.\;\ref{60099} hingegen nicht näher verwandt.
\pend
\pstart
Bereits in den Pariser Jahren (wohl um 1675/1676) hatte Leibniz die sieben Stoßregeln aus den \cite{00035}\textit{Principia} exzerpiert (\cite{01302}\textit{LSB} VI,~3 N.\;15, bes. S.\;216\,f.).
Allerdings beruhen die vorliegenden Auszüge auf einer späteren Lektüre, wie das Wasserzeichen bezeugt.
Sie sind mit sehr hoher Wahrscheinlichkeit mit 
Leibnizens \cite{01301}Marginalien im Hannoveraner Exemplar von 
Descartes' \cite{01300}\textit{Opera Philosophica}, Amsterdam 1650 
(GWLB, Leibn.~Marg.~6)
verwandt,
mit denen sie bedeutende inhaltliche, teilweise auch wörtliche Übereinstimmungen aufweisen.
Denn die Randbemerkungen im zweiten Teil der \textit{Principia}, §§~46\textendash52 (ediert in \cite{01301}\textit{LSB} VI, 4 N.\;335\textsubscript{1}) 
sind hauptsächlich einem Vergleich der cartesischen Stoßgesetze mit denjenigen gewidmet, die
\indexN{\textls{Malebranche}, Nicolas 1638\textendash1715}Malebranche
in den ersten vier Ausgaben von \cite{02080}\textit{De la Recherche de la vérité} aufgestellt hatte
(siehe \textsc{N.~Malebranche}, \cite{02080}\textit{Recherche}, l.~VI, chap.~9:
1.~Ausg. Paris 1674\textendash1675, tome II, S.\;493\textendash500 \lbrack483\textendash490\rbrack;
\cite{02079}\textit{MOC} 17,1, S.\;39\textendash44).
Die genannten Bemerkungen aus den Descartes-Marginalien finden sich an den einschlägigen Stellen von N.\;\ref{60099} wieder.
So z.\,B. die Kritik, Descartes habe einen der bedeutendsten Stoßfälle vernachlässigt 
(\cite{01301}\textit{LSB} VI, 4, S.\;1696 Anm.~1, hier \refpassage{35_10_12_038_8a}{35_10_12_038_8b}),
sowie die Feststellung von Malebranches Zustimmung (\glqq approbat\grqq) zu Descartes' ersten drei Regeln und zur fünften (dort Anm.~2 und 4, hier \refnote{FnA:35_10_12_038_approbet}).
Im Fall des Kommentars zu Malebranches Berichtigung der vierten Regel (hier \refpassage{35_10_12_038_5a}{35_10_12_038_7}) ist die Übereinstimmung mit der entsprechenden Randanmerkung (\glqq Malebranchius vult corpora simul ire post concursum, eadem qua ante motus quantitate\grqq, Anm.~3) sogar wörtlich.
Wegen dieser markanten Übereinstimmungen kann von einer engen Verwandtschaft beider Texte ausgegangen werden.
Des Weiteren erscheint die Annahme der Abhängigkeit der Auszüge N.\;\ref{60099} von den Marginalien (\cite{01301}\textit{LSB} VI, 4 N.\;335\textsubscript{1}) plausibel;
demzufolge hätte Leibniz bei der Exzerpierung der \cite{00035}\textit{Principia} die dort gemachten Notizen übernommen und in die gleichzeitig entstehenden Auszüge eingearbeitet.
\pend
\pstart
Die Marginalien wurden in \textit{LSB} VI, 4 auf den Zeitraum 1677\textendash1687 datiert (S.\;1695),
wobei die Datierung als unsicher erschlossen ausgewiesen ist.
Dies geschah mit Verweis auf Leibnizens intensive Beschäftigung mit Descartes zwischen 1683 und 1686
sowie auf die ausführliche Besprechung von dessen und Malebranches Stoßregeln 
im Rahmen der öffentlichen Auseinandersetzung von 1687 mit
\indexN{\textls{Catelan}, François, Abbé de gest. nach 1719}Catelan 
und Malebranche selbst
in den \cite{02002}\textit{Nouvelles de la République des Lettres}
Mit Hilfe der vorliegenden Auszüge N.\;\ref{60099} und mit Verweis auf andere Zeugnisse einer um 1685 stattfindenden Lektüre der \cite{02080}\textit{Recherche}
lässt sich die Descartes-Lektüre bezüglich der Stoßgesetze und die Entstehung der Marginalien näher eingrenzen.
\pend
\pstart
Bereits in der Pariser Zeit, wohl um 1675, hatte Leibniz Malebranches Bekanntschaft gemacht
(siehe \cite{02078}\textsc{Robinet} 1955, S.\;37\textendash40)
und die \cite{02080}\textit{Recherche} gelesen
(\cite{02078}ebd., bes. S.\;136), 
wobei nach dem aktuellen Kenntnisstand keine aus dieser Lektüre hervorgegangenen Texte überliefert sind.
Von einer zweiten, wie N.\;\ref{60099} in der Hannoveraner Zeit stattfindenden
Auseinandersetzung mit der \textit{Recherche} zeugen die
Anstreichungen und Marginalien in Leibnizens
\cite{02081}Handexemplar der 4. Ausgabe, Paris 1678\textendash1679
(\cite{02081}GWLB, Leibn.~Marg.~22)
sowie die ausführlichen Exzerpte daraus; 
die Texte sind in \cite{02082}\textit{LSB} VI, 4 N.\;348 geschlossen ediert.
Da dieses Exemplar zu den Büchern zählt, die laut einer 
von Leibniz im April 1685 der Hannoveraner Kammer vorgelegten
\cite{02083}Rechnung \glqq zu behuf der Fürstl. Bibliothec gekauffet worden\grqq\ (\textit{LSB} I,~4 N.\;148, S.\;190),
können seine \cite{02082}Anstreichungen und Anmerkungen frühestens zum Zeitpunkt des Erwerbs, etwa im Frühjahr 1685, entstanden sein 
(für die \cite{02082}Auszüge wird in \textit{LSB} VI, 4 die vermutliche Datierung \glqq 1686 bis 1699\,(?)\grqq\ vorgeschlagen).
\pend
\pstart
Leibniz hat sich bei der zweiten Lektüre nachweislich mit der einschlägigen Stelle der \textit{Recherche} beschäftigt
(siehe die Anstreichungen zu tome II, S.\;487: \cite{02082}\textit{LSB} VI,~4, S.\;1908 Z.~19\textendash23).
Es kann daher angenommen werden, dass erst die wohl um 1685 erfolgte zweite Malebranche-Lektüre Anlass und Stoff zu den kurzen, aber präzisen Angaben im Hannoveraner Descartes-Exemplar und in den vorliegenden Auszügen gab
\textendash\ obgleich diese Annahme angesichts der schon um 1675 erfolgten ersten Lektüre nicht zwingend notwendig ist.
Demzufolge würden 
der Kauf des Exemplars der \textit{Recherche} im Frühjahr 1685 einen Terminus post quem
und das Wasserzeichen einen Terminus ante quem (etwa bis Ende 1685)
für N.\;\ref{60099} abgeben.
Leibniz hätte also die Descartes-Auszüge etwa um 1685 verfasst,
innerhalb des bereits in \textit{LSB} VI,~4 (S.\;1695) konstatierten Zeitraums intensiverer Beschäftigung mit Descartes und im unmittelbaren Vorfeld der 
(am 6./16. Januar 1686 bei den \textit{Acta Eruditorum} eingesandten) 
\cite{01099}\textit{Brevis demonstratio erroris memorabilis Cartesii} (\textit{LSB} VI,~4 N.\;369).
\pend
\pstart
Zu erwähnen ist schließlich die inhaltliche und genetische Abhängigkeit einer weiteren Leibniz'schen Schrift  von N.\;\ref{60099}.
In der Polemik gegen 
\indexN{\textls{Catelan}, François, Abbé de gest. nach 1719}Catelan 
über das wahre Kraftmaß 
nimmt Leibniz die Stoßregeln 
\indexN{\textls{Descartes} (Cartesius, des Cartes), René 1596\textendash1650}Descartes' und \indexN{\textls{Malebranche}, Nicolas 1638\textendash1715}Malebranches 
öffentlich ins Visier.
In sein erstes Schreiben
(\cite{02086}\glqq Replique de M.\,L. à M. l'Abbé D.\,C.\grqq, in \textit{NRL} Februar 1687, S.\;131\textendash145)
hat er einige Ergebnisse und Textteile von N.\;\ref{60099} einfließen lassen.
So ist eine exakte Übernahme des hier aufgestellten Gegenbeispiels zur dritten cartesischen Regel, einschließlich der von Leibniz erst im Laufe der Bearbeitung festgelegten Zahlenwerte, festzustellen (vgl. S.\;137\,f. der \glqq Replique\grqq\ mit \refpassage{35_10_12_038_6a}{35_10_12_038_6b}).
Auch der Schlusskommentar über die absurde Folge dieser Regel (d.\,h. den Verlust beinahe der Hälfte der Kraft) findet sich dort genau wieder, jedoch auf Französisch.
Die genetische Abhängigkeit dieser Textteile von den vorliegenden Auszügen 
wird weiter durch den Umstand bestätigt, dass 
eine unmittelbar benachbarte Signatur (LH~XXXV~10, 12~Bl.~39\textendash40)
eins der erhaltenen Konzepte der "Replique" überliefert,
das als Vorlage für die zitierte Stelle gedient hat und die Übernahme des Beispiels und der Berechnungen aus N.\;\ref{60099} bereits aufweist.
\pend
\end{ledgroup}
\newpage
\selectlanguage{latin}
\vspace{8mm}
\pstart%
\centering
\href{\Scan{LH_35_10_12_038r}}%
{\lbrack 38~r\textsuperscript{o}\rbrack}
\hfill
Notae in 
\edtext{septem}{%
\lemma{}%
\Bfootnote{%
septem \textit{erg.~L}}} 
regulas motus\indexS{regulae motus Cartesii} 
\indexN{\textls{Descartes} (Cartesius, des Cartes), René 1596\textendash1650}Cartesii
\hfill
\phantom{\lbrack 38~r\textsuperscript{o}\rbrack}
\pend
\begin{edfigure}[1.0em][1.0em]    	%
\includegraphics[scale=1]{%
gesamttex/edit_VIII-4/images/LH_35_10_12_038_d1_038r.pdf}
\edcaption{%
\edtext{\lbrack\textit{Fig.~1}\rbrack}{%
\lemma{\lbrack\textit{Fig.~1}\rbrack}%
\Cfootnote[nonum]{%
Siehe \indexN{\textls{Descartes} (Cartesius, des Cartes), René 1596\textendash1650}\textsc{R.~Descartes},
\cite{00035}\textit{Principia Philosophiae}, Pars~II, §~46, Amsterdam 1644, 
S.\;60 (\textit{DO} VIII.1, S.\;68).}}}
\end{edfigure}
\pstart\noindent
\textls{Regula prima} est:\enskip 
\edtext{%
\textit{Si duo corpora \textit{B} et \textit{C} essent plane aequalia, et aeque velociter}
\edtext{\textit{moverentur} in contrarias partes,\indexS{corpora aequalia in contrarias partes mota}}{%
\lemma{\textit{moverentur}}%
\Bfootnote{%
\textit{(1)}~ita ut occurrere %
\textit{(2)}~in contrarias partes,~\textit{L}}}
\edtext{\textit{\textit{B} a dextra ad sinistram}, \textit{C} autem contra\lbrack,\rbrack}{%
\lemma{}%
\Bfootnote{%
\textit{B} \lbrack...\rbrack\ contra \textit{erg.~L}}} 
reflectentur post concursum\indexS{concursus} eadem qua venerunt celeritate.%
}{\lemma{\textit{Si duo} \lbrack...\rbrack\ celeritate}%
\Cfootnote{%
\cite{00035}a.\,a.\,O., §~46, S.\;59\,f. mit Auslassungen.}}	%
Hanc regulam solam ex septem istis\indexS{regulae motus Cartesii} admitto.
\pend
\pstart
\textls{Regula secunda}:\enskip 
\edtext{%
\textit{Si \textit{B} esset}
\edtext{\textit{tantillo} majus\indexS{corpus majus}}{%
\lemma{\textit{tantillo}}%
\Bfootnote{%
\textit{(1)}~\textit{major} %
\textit{(2)}~majus~\textit{L}}}
\textit{quam \textit{C}, caeteris positis ut prius, tunc solum \textit{C} reflecteretur 
et utrumque versus sinistram eadem celeritate moveretur.}%
}{\lemma{\textit{Si \textit{B}} \lbrack...\rbrack\ \textit{moveretur}}%
\Cfootnote{%
\cite{00035}a.\,a.\,O., §~47, S.\;60.}}	%
Haec regula\indexS{regulae motus Cartesii} rationi contraria 
\edtext{est, tantillo}{%
\lemma{est,}%
\Bfootnote{%
\textit{(1)}~scilicet %
\textit{(2)}~usque a %
\textit{(3)}~tantillo~\textit{L}}}
scilicet adjecto ad \textit{B}, quod minus sit data quantitate, effectum plane contrarium nasci, 
ut cum antea iret ad dextram, 
\edtext{nunc eadem}{%
\lemma{nunc}%
\Bfootnote{%
\textit{(1)}~contra %
\textit{(2)}~eadem~\textit{L}}}
celeritate in contrarias partes tendat\lbrack,\rbrack\ 
cum tamen insensibilis mutatio in hypothesi,\indexS{mutatio insensibilis in hypothesis}
debeat etiam in consequentiis insensibilem mutationem\indexS{mutatio insensibilis in consequentiis}
\edtext{efficere. Praeterea in hoc casu, quo \textit{B} est tantillo majus\indexS{corpus majus} quam \textit{C}\lbrack,\rbrack\ directio ipsius \textit{B} sine ulla refractione\indexS{refractio} seu diminutione vincit}{%
\lemma{efficere.}%
\Bfootnote{%
\textit{(1)}~\textls{Regula tertia} %
\textit{(2)}~Ut taceam %
\textit{(3)}~Praeterea %
\textit{(a)}~hoc modo %
\textit{(b)}~sola directio ipsius \textit{B} vincit %
\textit{(c)}~in hoc \lbrack...\rbrack\ ipsius \textit{B} 
\bars{sine ulla refractione seu diminutione \textit{erg.}} vincit~\textit{L}}}
nulla ratione habita directionis ipsius \textit{C}, 
cum tamen semper utraque in nova directione resultante aliquid efficere debeat, 
et id quod vincit semper etiam nonnihil de suo remittat. 
\pend
\pstart
\textls{Regula tertia}:\enskip 
\edtext{%
\textit{Si mole essent aequalia}\indexS{corpora aequalia} et \textit{B}
\textit{tantillo celerius moveretur quam \textit{C}}\lbrack,\rbrack\indexS{corpus celerius}
\textit{non} 
\edtext{\textit{tantum} post occursum\indexS{occursus} \textit{ambo} moverentur \textit{versus sinistram},}{%
\lemma{\textit{tantum}}%
\Bfootnote{%
\textit{(1)}~\textit{ambo} %
\textit{(2)}~post occursum \textit{ambo} %
\textit{(a)}~\textit{pergerent moveri versus sinistram} %
\textit{(b)}~moverentur \textit{versus sinistram},~\textit{L}}}
\ledpb%
\edlabel{35_10_12_038_2a}%
\edtext{}{%
\xxref{35_10_12_038_2a}{35_10_12_038_2b}%
\lemma{}%
\Afootnote{%
\textit{Am Rand}:
\hspace*{2em}%
secundum \indexN{\textls{Descartes} (Cartesius, des Cartes), René 1596\textendash1650}Cartesium
\vspace{1ex}
\newline
\hspace*{2em}\protect\begin{tabular}{lll}
ante conflictum:\indexS{conflictus} & $6B+4C=m$\textsuperscript{\lbrack a\rbrack} & $B=C$\\
post conflictum\lbrack:\rbrack\indexS{conflictus} & $5B+5C=(m)$ & $m=(m)$\textsuperscript{\lbrack b\rbrack}
\\[1.0ex]
& $36B+16C=v$ & $B=C$\\
& $25B+25C=(v)$ &
\protect\end{tabular}
\vspace{1ex}
\newline
Ergo \textit{v} non est $=(v)$ nam \lbrack52\rbrack\textsuperscript{\lbrack c\rbrack} non est $=50$
\\[1.5ex] %
{\footnotesize 
\textsuperscript{\lbrack a\rbrack} $6B+4C=m$: Leibniz hat in der Randanmerkung die zwei Körper zunächst durchgehend mit \textit{A} und \textit{B} bezeichnet und die Bezeichnungen nachträglich der Vorlage angepasst.\quad
\textsuperscript{\lbrack b\rbrack} $m=(m)$ \bars{secundum me \textit{gestr.}} $36B+16C=v$ \textit{L}\quad
\textsuperscript{\lbrack c\rbrack} 46 \textit{L ändert Hrsg.}}}}%
\textit{sed etiam transferretur ex \textit{B} in \textit{C} media pars celeritatis}
\edtext{\textit{quae hoc}}{%
\lemma{\textit{quae}}%
\Bfootnote{%
\textit{(1)}~hac %
\textit{(2)}~\textit{hoc}~\textit{L}}}
\textit{ab illo excederetur}, ita ut postea amborum aequalis sit celeritas.%
}{\lemma{\textit{Si mole} \lbrack...\rbrack\ celeritas}%
\Cfootnote{%
\cite{00035}a.\,a.\,O., §~48, S.\;60.}}	%
Haec regula\indexS{regulae motus Cartesii} est contra naturam percussionis corporum durorum,\indexS{corpus durum}\indexS{percussio corporum durorum} 
semper enim duo corpora 
\edtext{dura post}{%
\lemma{dura}%
\Bfootnote{%
\textit{(1)}~quae sese %
\textit{(2)}~post~\textit{L}}}
conflictum\indexS{conflictus} a se invicem dissiliunt, nec unquam simul procedunt.%
\edlabel{35_10_12_038_2b}
Mollia\indexS{corpus molle} non dissiliunt, quia pars virium\indexS{pars virium} perditur in conflictu,\indexS{conflictus} 
absumiturque a mollium partibus insensibilibus\indexS{partes molles insensibiles} 
quae inde motum intestinum inobservabilem\indexS{motus intestinus inobservabilis} accipiunt. 
Sed quod est 
\edlabel{35_10_12_038_3a}%
\edtext{}{%
\xxref{35_10_12_038_3a}{35_10_12_038_3b}%
\lemma{multo}%
\Bfootnote{%
\textit{(1)}~absurdius ostendam effectum hoc modo fore potentiorem causa, ac proinde haberi motum perpetuum Mechanicum. %
\textit{(2)}~minus %
\textit{(3)}~absurdius, ostendam %
\textit{(a)}~caus %
\textit{(b)}~plurimum potentiae perdi sine ulla ratione. %
\textit{(aa)}~Ex plano inclinato %
\textit{(bb)}~In planis inclinatis %
\textit{(aaa)}~\textit{LS} %
\textit{(bbb)}~\textit{LP}, et
\textit{(aaaa)}~\textit{NQ} %
\textit{(bbbb)}~\textit{NV} %
\textit{(cccc)}~\textit{NMQ}~\textit{L}}}%
multo absurdius, ostendam plurimum potentiae\indexS{potentia} perdi\indexS{amissio potentiae} sine ulla ratione.
\pend
\begin{edfigure}%
\includegraphics[scale=1]{%
gesamttex/edit_VIII-4/images/LH_35_10_12_038_d2_038r.pdf}
\edcaption[0.25em]{\lbrack\textit{Fig.~2}\rbrack}
\end{edfigure}
\pstart
\edlabel{35_10_12_038_6a}%
In planis inclinatis\indexS{planum inclinatum} \textit{LP}, et \textit{NMQ}%
\edlabel{35_10_12_038_3b}
in planum subjectum
\edtext{horizontale\indexS{planum horizontale} \textit{PQ} descendant globi per omnia aequales\indexS{globi aequales} \textit{B} et \textit{C}, ille \textit{B} ex altitudine \textit{LS} quae sit 10000 pollicum, hic}{%
\lemma{horizontale}%
\Bfootnote{%
\textit{(1)}~\textit{MP}
\textit{(2)}~\textit{PQ}
\textit{(a)}~corpora aequalia \textit{B} et \textit{C} %
\textit{(b)}~globi per omnia aequales %
\textit{(aa)}~et ex %
\textit{(bb)}~\textit{B} et \textit{C}, ille \bars{\textit{B} \textit{erg.}} ex altitudine
\textit{(aaa)}~\textit{NT}, 
\textit{(aaaa)}~quae sit
\textit{(bbbb)}~quae sit 3
\textit{(bbb)}~quae sit %
\textit{(ccc)}~\textit{LS} quae sit %
\textit{(aaaa)}~36, haec ex alt %
\textit{(bbbb)}~36 %
\textit{(aaaaa)}~pollicum, haec ex altitudine %
\textit{(bbbbb)}~pedum, %
\textit{(ccccc)}~10000 pollicum, %
\textit{(aaaaa-a)}~haec %
\textit{(bbbbb-b)}~hic~\textit{L}}}
\textit{C} ex altitudine
\edtext{\lbrack\textit{MT}\rbrack}{%
\lemma{}%
\Bfootnote{%
\textit{NT} \textit{L ändert Hrsg.}}}
quae 
\edtext{sit unius pollicis, occurrant sibi in plano horizontali\indexS{planum horizontale} \textit{PQ}, nempe \textit{B} celeritate ut 100, et \textit{C} celeritate ut 1. Post concursum\indexS{concursus} autem in \textit{\ind{2}B\ind{2}C}, ambo}{%
\lemma{sit}%
\Bfootnote{%
\textit{(1)}~16 \bars{pedum \textit{streicht Hrsg.}}  %
\textit{(2)}~unius pollicis, %
\textit{(a)}~sibique occur %
\textit{(b)}~si %
\textit{(c)}~occurrent %
\textit{(d)}~occurrant sibi in plano
\textit{(aa)}~celeritatibus, %
\textit{(bb)}~\textit{B} celeritat %
\textit{(cc)}~horizontali %
\textit{(aaa)}~, \textit{B} haben %
\textit{(bbb)}~\textit{PQ}, nempe \textit{B} celeritate ut
\textit{(aaaa)}~6 %
\textit{(bbbb)}~100, et \textit{C} celeritate ut
\textit{(aaaaa)}~4 %
\textit{(bbbbb)}~1. Post concursum
\bars{autem \textit{erg.}} %
\textit{(aaaaa-a)}~\textit{\ind{2}C\ind{2}B} %
\textit{(bbbbb-b)}~in %
\textit{(aaaaa-aa)}~\textit{\ind{2}B} %
\textit{(bbbbb-bb)}~\textit{\ind{2}B\ind{2}C} %
\textit{(ccccc-c)}~in \textit{\ind{2}B\ind{2}C}, %
\textit{(aaaaa-aa)}~ipsum %
\textit{(bbbbb-bb)}~ambo~\textit{L}}}
si regulae 
\edtext{credimus inde simul tendent versus \textit{Q} celeritate, ut 
$\displaystyle\frac{101}{2}$}{%
\lemma{credimus}%
\Bfootnote{%
\textit{(1)}~\textit{B} secum abripiens \textit{C} assurget cum ipso ad \textit{N}, 
\textit{(2)}~ambo ibunt inde %
\textit{(3)}~\bars{ambo \textit{gestr.}} inde simul %
\textit{(a)}~ibunt %
\textit{(b)}~tendent versus
\textit{(aa)}~\textit{N} %
\textit{(bb)}~\textit{Q} celeritate, ut
\textit{(aaa)}~5, atque %
\textit{(bbb)}~\protect\rule[-0ex]{0pt}{3.5ex}$\displaystyle\frac{101}{2}$~\textit{L}}}
\rule[-0ex]{0pt}{3.5ex}inde porro assurgent usque ad \textit{N} posito 
\edtext{altitudinem \textit{NV} essent $\displaystyle\frac{10201}{4}$.}{%
\lemma{altitudinem}%
\Bfootnote{%
\textit{(1)}~\textit{M} %
\textit{(2)}~\textit{NV} %
\textit{(a)}~esse ut 25. %
\textit{(b)}~essent \protect\rule[-2ex]{0pt}{0ex}$\displaystyle\frac{10201}{4}$.~\textit{L}}}
\pend
\begin{edfigure}[1em]%
\includegraphics[scale=1]{%
gesamttex/edit_VIII-4/images/LH_35_10_12_038_d3_038r.pdf}
\edcaption[0.25em]{\lbrack\textit{Fig.~3}\rbrack}
\end{edfigure}
\pstart
Cum ergo duo corpora 
\edtext{aequalia quae ponamus esse singula unius librae}{%
\lemma{aequalia}%
\Bfootnote{%
\textit{(1)}~simul
\textit{(1)}~quae \lbrack...\rbrack\ librae~\textit{L}}}
initio elevata\indexS{elevatio unius librae} fuissent 
\edtext{unum ad altitudinem ut 10000, alterum ad altitudinem ut 1,}{%
\lemma{unum ad altitudinem ut}%
\Bfootnote{%
\textit{(1)}~25 %
\textit{(2)}~10000, alterum ad altitudinem ut
\textit{(a)}~25 %
\textit{(b)}~46 %
\textit{(c)}~1,~\textit{L}}}
quod est idem ac 
\edtext{si una}{%
\lemma{si}%
\Bfootnote{%
\textit{(1)}~unum %
\textit{(2)}~una~\textit{L}}}
libra esset 
\edtext{elevata\indexS{elevatio unius librae} ad pollices 10001,}{%
\lemma{elevata ad}%
\Bfootnote{%
\textit{(1)}~pedes 46, %
\textit{(2)}~pollices 10001,~\textit{L}}}
nunc vero duae librae sint 
\edtext{elevatae ad pollices $10201:4$}{%
\lemma{elevatae ad}%
\Bfootnote{%
\textit{(1)}~\bars{pedes \textit{streicht Hrsg.}} 25 %
\textit{(2)}~pollices $10201:4$~\textit{L}}}
quod idem est ac unam libram esse 
\edtext{elevatam\indexS{elevatio unius librae} ad pollices $10201:2$}{%
\lemma{elevatam}%
\Bfootnote{%
\textit{(1)}~ad %
\textit{(a)}~pedes 50, %
\textit{(b)}~pollices $10201:2$, patet effectum seu statum posteriorem esse potentiorem 
causa seu statu priore, adeoque haberi posse motum perpetuum 
\bars{Mechanicum \textit{gestr.}}, quod est absurdum. %
\textit{(2)}~ad pollices $10201:2$~\textit{L}}}
seu 5100 cum dimidio, cum 
\edtext{antea perinde}{%
\lemma{antea}%
\Bfootnote{%
\textit{(1)}~vir %
\textit{(2)}~tantum virium esset, quantum elevare posset libram ad %
\textit{(3)}~libra elevat %
\textit{(4)}~perinde~\textit{L}}}
esset ac si libra esset elevata\indexS{elevatio unius librae} ad pollices 10001. Itaque dimidia 
\edtext{fere}{%
\lemma{}%
\Bfootnote{fere \textit{erg.~L}}} 
pars potentiae\indexS{potentia} periit,\indexS{amissio partis potentiae} 
nulli rei impensa, ac proinde sine ulla causa, quod absurdum est.%
\edlabel{35_10_12_038_6b}	%
Notabile est, quo majus discrimen sumitur inter aequalium corporum\indexS{corpora aequalia} celeritates, 
eo major apparet error ita ut tandem assumto 
\edtext{infinito, seu posita quiete, alterutrius dimidium potentiae amittatur.\indexS{amissio dimidiae potentiae}}{%
\lemma{infinito}%
\Bfootnote{%
\textit{(1)}~error %
\textit{(2)}~, seu ratio virium %
\textit{(3)}~seu posita quiete, alterutrius %
\textit{(a)}~ratio virium fiat 3 %
\textit{(b)}~tres quartae potentiae amittantur
\textit{(c)}~dimidium potentiae amittatur.~\textit{L}}}
\pend
\pstart
Si vero aequalitas sit 
\edtext{in celeritatibus}{%
\lemma{in}%
\Bfootnote{%
\textit{(1)}~rationibus %
\textit{(2)}~celeritatibus~\textit{L}}}
error erit minimus, id est nullus, 
\href{\Scan{LH_35_10_12_038v}}%
{\lbrack 38~v\textsuperscript{o}\rbrack}\ seu nihil virium amittetur,\indexS{amissio virium} quod calculo ostendere placet. Sint celeritates \textit{a} 
\edtext{et \textit{b}; corporibus}{%
\lemma{et \textit{b};}%
\Bfootnote{%
\textit{(1)}~erit vi %
\textit{(2)}~corporibus~\textit{L}}}
igitur existentibus aequalibus, ac proinde inter calculandum omissis\lbrack,\rbrack\ erunt vires
$a^2+b^2$. ante ictum,\indexS{ictus}\indexS{vis ante ictum} at post ictum\indexS{ictus}
\edtext{secundum regulam}{%
\lemma{secundum}%
\Bfootnote{%
\textit{(1)}~celeritatem %
\textit{(2)}~regulam~\textit{L}}}
\indexN{\textls{Descartes} (Cartesius, des Cartes), René 1596\textendash1650}Cartesii\indexS{regulae motus Cartesii} celeritas 
\edtext{amborum}{%
\lemma{}%
\Bfootnote{amborum \textit{erg.~L}}} 
erit 
\edtext{$a+b,:2$. cujus quadratum\indexS{quadratum celeritatis}}{%
\lemma{$a+b,:2$.}%
\Bfootnote{%
\textit{(1)}~unde vi %
\textit{(2)}~cujus quadratum~\textit{L}}}
est $a^2+b^2+2ab,:4$ 
quod ductum\indexS{quadratum celeritatis ductum in molem} in 2 molem\indexS{moles} amborum dabit $a^2+bb+2ab,:2$, 
quo detracto a prioribus viribus,\indexS{vis} restabit: $2a^2+2b^2-a^2-b^2-2ab,:2=a^2+b^2-2ab,:2$. differentia 
\edtext{virium,\indexS{vis} cujus}{%
\lemma{virium,}%
\Bfootnote{%
\textit{(1)}~qua de ha %
\textit{(2)}~cujus~\textit{L}}}
ratio ad vires\indexS{vis} priores est $a^2+b^2-2ab,:\llcorner 2\;\overline{a^2+b^2}$. Seu pro 
$a:b$ ponendo \textit{y}, fiet 
\edtext{$yy+1-2y,:2yy+2$ quod ponatur $=z$. Et quaerendo}{%
\lemma{$yy+1-2y,:2yy+2$}%
\Bfootnote{%
\textit{(1)}~quod si ponatur $=z$ et quaeratur %
\textit{(2)}~quod ponatur $=z$
\textit{(a)}~cujus si %
\textit{(b)}~. Et quaerendo~\textit{L}}}
maximum et minimum\indexS{maximum et minimum} 
\edtext{fiet %
\edtext{mea methodo\indexS{methodus mea}}{%
\lemma{mea methodo\indexS{methodus mea}}%
\Cfootnote{%
Bereits um 1675\textendash1676 hatte Leibniz Fragen \glqq de maximis et minimis\grqq\ anhand seiner neuen Differentialmethode beantwortet (siehe die Schriften zur Infinitesimalmathematik in \textit{LSB} VII, 5 \textit{passim}); seine erste einschlägige Veröffentlichung ist die \cite{02077}\textit{Nova methodus pro maximis et minimis} (\cite{01023}\textit{AE} Oktober 1684, S.\;467\textendash473).}}
statim: $\overline{2y-2}\,\cdot\,\overline{2yy+2}=\overline{yy+1-2y}\,\cdot\,4y$,}{%
\lemma{fiet}%
\Bfootnote{%
\textit{(1)}~: $\overline{2y-2}\cdot\overline{2yy+2}=\overline{yy+1-2y}\cdot 4y$ ex mea methodo de maximis et minimis, %
\textit{(2)}~mea methodo statim: $\overline{2y-2}\cdot\overline{2yy+2}=\overline{yy+1-2y}\cdot 4y$,~\textit{L}}}
itaque fiet: $4y^3-4yy+4y-4=4y^3-8yy+4y$. seu $4yy=4$ seu $y=1$. Quod si ponamus 
\edtext{$y=0$. si scilicet \textit{a} infinitae parvitatis respectu \textit{b}.}{%
\lemma{$y=0$. si}%
\Bfootnote{%
\textit{(1)}~ratio %
\textit{(2)}~scilicet %
\textit{(a)}~$a:b$ sit %
\textit{(b)}~rat %
\textit{(c)}~\textit{a} infinitae parvitatis respectu \textit{b}.~\textit{L}}}
\edlabel{35_10_12_038_1a}%
\edtext{}{\xxref{35_10_12_038_1a}{35_10_12_038_1b}%
\lemma{}%
\Bfootnote{%
fiet $z=1:2$. 
\bars{Malebranchius \textit{erg. u. gestr.}}
Regula~\textit{L}}}%
fiet $z=1:2$.
\pend
\pstart
Regula\edlabel{35_10_12_038_1b} IVta.
\edtext{%
\textit{Si \textit{C} quiesceret et esset paulo majus\indexS{corpus majus quiescens} 
quam \textit{B}, nunquam \textit{C} moveretur} 
a \textit{B} sed ipsum \textit{repelleret in contrariam partem} ea vi qua venit.%
}{\lemma{\textit{Si} \lbrack...\rbrack\ venit}%
\Cfootnote{%
\cite{00035}a.\,a.\,O., §~49, S.\;60 mit Auslassungen.}}	%
Ita \indexN{\textls{Descartes} (Cartesius, des Cartes), René 1596\textendash1650}Cartesius, 
eoque casu fateor eandem potentiam mansuram,\indexS{conservatio potentiae} 
quia eadem in eodem corpore celeritas manet. Sed corpus 
\edtext{quiescens\indexS{corpus quiescens} a minori}{%
\lemma{quiescens}%
\Bfootnote{%
\textit{(1)}~nulla vi %
\textit{(2)}~ab aequali %
\textit{(3)}~a minori~\textit{L}}}
nulla vi impelli posse, plane incredibile est\lbrack,\rbrack\
\edtext{cum tamen ipse 
\edtext{velit ab aequali}{%
\lemma{velit}%
\Bfootnote{%
\textit{(1)}~a maj %
\textit{(2)}~ab aequali~\textit{L}}}
posse impelli.}{\lemma{cum \lbrack...\rbrack\  impelli}%
\Cfootnote{%
Siehe \indexN{\textls{Descartes} (Cartesius, des Cartes), René 1596\textendash1650}Descartes' sechste Regel, 
\cite{00035}a.\,a.\,O., §~51, S.\;61.}}	%
\edtext{Itaque utcunque}{%
\lemma{Itaque}%
\Bfootnote{%
\textit{(1)}~vel %
\textit{(2)}~utcunque~\textit{L}}}
exigua accessione infinitam acquiret Resistentiam,\indexS{Resistentia} quod absurdum est. Quodsi volebat 
\edtext{quietem tantam vim\indexS{vis}}{%
\lemma{quietem}%
\Bfootnote{%
\textit{(1)}~, et vim %
\textit{(2)}~corpus quiescens %
\textit{(3)}~tantam vim~\textit{L}}}
ad 
\edtext{resistend\lbrack u\rbrack m}{%
\lemma{}%
\Bfootnote{resistendam \textit{L ändert Hrsg.}}} 
habere, debebat asserere 
\edtext{etiam minus}{%
\lemma{etiam}%
\Bfootnote{%
\textit{(1)}~majus %
\textit{(2)}~minus~\textit{L}}}
corpus quiescens\indexS{corpus quiescens} ab aequali aut etiam majori non posse impelli nisi quod esset infinities majus.
\edlabel{35_10_12_038_5a}%
\edtext{}{%
\xxref{35_10_12_038_5a}{35_10_12_038_5b}%
\lemma{Malebranchius \lbrack...\rbrack\ omnibus}%
\Cfootnote{%
\indexN{\textls{Malebranche}, Nicolas 1638\textendash1715}\textsc{N.~Malebranche}, \cite{02081}\textit{Recherche}, l.~VI, p.~2, chap.~9, Paris 1678, S.\;486\textendash492, hier S.\;487\,f.
(\cite{02079}\textit{MOC} 17, 1, S.\;39\textendash44, hier S.\;40).}}%
Malebranchius\indexN{\textls{Malebranche}, Nicolas 1638\textendash1715}
vult in casu regulae 4tae\indexS{regulae motus Cartesii} corpora simul ire post concursum,\indexS{concursus} 
eadem qua ante motus quantitate,\indexS{quantitas motus}\edlabel{35_10_12_038_7} %
idemque statuit in casibus caeteris 
\edtext{omnibus,\edlabel{35_10_12_038_5b}
sed ita aequalitas causae et}{%
\lemma{omnibus,}%
\Bfootnote{%
\textit{(1)}~sed hoc ego plane absurdum %
\textit{(2)}~sed %
\textit{(a)}~ho %
\textit{(b)}~ita %
\textit{(aa)}~rursus %
\textit{(bb)}~aequalitas 
\textit{(aaa)}~motus %
\textit{(bbb)}~causae et~\textit{L}}}
effectus\indexS{aequalitas causae et effectus} violabitur, 
ut ad modum regulae 3\indexS{regulae motus Cartesii} calculando facile cognoscere licebit. 
\pend
\pstart
Regula V. 
\edtext{%
\edtext{Si corpus}{%
\lemma{Si}%
\Bfootnote{%
\textit{(1)}~\textit{C} minus, %
\textit{(2)}~corpus~\textit{L}}}
minus\indexS{corpus minus quiescens} \textit{C},
\edtext{quiescit\lbrack,\rbrack\ ambo post concursum\indexS{concursus}}{%
\lemma{quiescit}%
\Bfootnote{%
\textit{(1)}~moveretur %
\textit{(2)}~ambo post concursum~\textit{L}}}
simul ibunt priore quantitate motus.\indexS{quantitas motus}%
}{\lemma{Si corpus \lbrack...\rbrack\ motus}%
\Cfootnote{%
\indexN{\textls{Descartes} (Cartesius, des Cartes), René 1596\textendash1650}\textsc{Descartes},
\cite{00035}\textit{Principia Philosophiae}, §~50, S.\;60\,f.}}	%
Haec regula\indexS{regulae motus Cartesii} non minus 
\edtext{quam tertia}{%
\lemma{quam}%
\Bfootnote{%
\textit{(1)}~quarta %
\textit{(2)}~tertia~\textit{L}}}
contra aequalitatem causae et Effectus\indexS{aequalitas causae et effectus} peccat\lbrack,\rbrack\ 
quamvis 
\edtext{eam pariter ac primam secundam et 3tiam 
\indexN{\textls{Descartes} (Cartesius, des Cartes), René 1596\textendash1650}Cartesii 
regulam\indexS{regulae motus Cartesii}
Malebranchius\indexN{\textls{Malebranche}, Nicolas 1638\textendash1715} approbet.}{%
\lemma{eam \lbrack...\rbrack\ approbet}%
\applabel{FnA:35_10_12_038_approbet}%
\Cfootnote{%
\indexN{\textls{Malebranche}, Nicolas 1638\textendash1715}\textsc{Malebranche}, \cite{02081}\textit{Recherche}, a.\,a.\,O., S.\;487
(\cite{02079}\textit{MOC} 17, 1, S.\;40).}} 
\pend
\pstart
Reg. VI.
\edlabel{35_10_12_038_4a}%
\edtext{}{%
\xxref{35_10_12_038_4a}{35_10_12_038_4b}%
\lemma{\textit{Si corpus} \lbrack...\rbrack\ quarta}%
\Cfootnote{%
\indexN{\textls{Descartes} (Cartesius, des Cartes), René 1596\textendash1650}\textsc{Descartes},
\cite{00035}\textit{Principia Philosophiae}, §~51, S.\;61 mit Auslassungen.}}%
\edtext{%
\textit{Si corpus quiescens\indexS{corpus quiescens}}
\edtext{\textit{esset aequale}}{%
\lemma{}%
\Bfootnote{%
\textit{esset} 
\bars{\textit{accuratissime} \textit{gestr.}}
\textit{aequale}~\textit{L}}}
\textit{corpori \textit{B} versus illud moto, \textit{B} reflecteretur cum tribus}
quartis velocitatis,}{%
\lemma{}%
\Afootnote{%
\textit{Am Rand:} 
\protect\begin{tabular}[t]{l}
\;\,1 \\[1.5ex]
$\displaystyle\frac{9}{16}+\displaystyle\frac{1}{16}$
\protect\end{tabular}}}
\textit{C} vero impelleretur cum una 
\edtext{quarta.\edlabel{35_10_12_038_4b} Non}{%
\lemma{quarta.}%
\Bfootnote{%
\textit{(1)}~Hic vellem exurgeret aliquis, qui vel ex ipsis Cartesii principiis rationem redderet %
\textit{(a)}~divi %
\textit{(b)}~\bars{talium \textit{streicht Hrsg.}} numerorum; %
\textit{(2)}~Non~\textit{L}}}
possum hic divinare, unde divisio velocitatum in quatuor partes, et 
\edtext{cur inter}{%
\lemma{cur}%
\Bfootnote{%
\textit{(1)}~corpus %
\textit{(2)}~inter~\textit{L}}}
duo aequalia\indexS{corpora aequalia} ita debeat distribui velocitas, ut tripla sit ejus quod reflectitur. 
\edtext{Caeterum si}{%
\lemma{Caeterum}%
\Bfootnote{%
\textit{(1)}~secundum Cartesii Calculum cum prio %
\textit{(2)}~si~\textit{L}}}
haec regula\indexS{regulae motus Cartesii} vera est, cum prius adesset potentia\indexS{potentia} elevandi v.\,g. libram\indexS{elevatio unius librae} 
\edtext{ad pedes 16,}{%
\lemma{ad}%
\Bfootnote{%
\textit{(1)}~unum pedem, %
\textit{(2)}~pedes 16,~\textit{L}}}
nunc post concursum tantum potentia\indexS{potentia} adest elevandi libram\indexS{elevatio unius librae} ad pedes 10, et 
\edtext{ita tertia}{%
\lemma{ita}%
\Bfootnote{%
\textit{(1)}~dimidia %
\textit{(2)}~tertia~\textit{L}}}
pars virium\indexS{vis} sine causa periit. 
\pend
\pstart
Reg. VII. 
\edtext{%
Si \textit{B} insequeretur \textit{C} atque attingeret, essetque \textit{B} minus, 
\textit{sed excessus celeritatis\indexS{excessus celeritatis} in \textit{B} major quam excessus magnitudinis\indexS{excessus magnitudinis} in \textit{C}, ambo}
\edtext{post concursum\indexS{concursus}}{%
\lemma{post}%
\Bfootnote{%
\textit{(1)}~ictum %
\textit{(2)}~concursum~\textit{L}}}
\textit{aeque celeriter moverentur} servata eadem quantitate motus.\indexS{quantitas motus}\indexS{conservatio quantitatis motus}
Sin \textit{minor esset excessus celeritatis\indexS{excessus celeritatis} quam magnitudinis,\indexS{excessus magnitudinis} 
\textit{B} in contrariam partem reflecteretur, et motum\indexS{motus} suum omnem retineret.}%
}{\lemma{Si \textit{B} \lbrack...\rbrack\ \textit{retineret}}%
\Cfootnote{%
\cite{00035}a.\,a.\,O., §~52, S.\;61\,f. mit Auslassungen.}}	%
Itaque et \textit{C} pergeret ea qua prius celeritate. 
Haec regula\indexS{regulae motus Cartesii} quam perturbata sit indicatu facile 
\edtext{est, nec}{%
\lemma{est,}%
\Bfootnote{%
\textit{(1)}~quis credat %
\textit{(2)}~nec~\textit{L}}}
facile apparet quomodo haec distinctio prioribus regulis\indexS{regulae motus Cartesii} concilietur, 
ubi compensatae per celeritatem\indexS{celeritas} magnitudinis\indexS{magnitudo} nullam habuerat rationem. 
Quemadmodum hic quoque eum omittit casum, 
\edtext{quo praecise reciproca}{%
\lemma{quo}%
\Bfootnote{%
\textit{(1)}~praecisa
\textit{(2)}~praecise
\textit{(a)}~ea ratio est %
\textit{(b)}~reciproca~\textit{L}}}
est magnitudinum et celeritatum 
\edtext{ratio.\indexS{ratio magnitudinum et celeritatum reciproca} Sed et in priori casu fit ut effectus sit  inaequalis causae.\indexS{effectus inaequalis causae}}{%
\lemma{ratio.}%
\Bfootnote{%
\textit{(1)}~Et vero quis credat corpus in aliquam partem tendens, %
\textit{(a)}~alteri %
\textit{(b)}~(in casu posteriori regulae) alteri eodem tendenti atque ipsum insequenti %
\textit{(2)}~Sed %
\textit{(a)}~in priori %
\textit{(b)}~et in priori casu %
\textit{(aa)}~oritur %
\textit{(bb)}~fit ut effectus sit
\textit{(aaa)}~potentior \bars{causa \textit{streicht Hrsg.}} %
\textit{(bbb)}~inaequalis causae.~\textit{L}}}
\edtext{Nam sit \textit{B},}{%
\lemma{Nam}%
\Bfootnote{%
\textit{(1)}~sit \textit{B} 1, et \textit{C} sit 2  %
\textit{(2)}~sit \textit{B} 1 celeritate 1, et \textit{C} 2, celeritate 1 %
\textit{(3)}~sit \textit{B},~\textit{L}}}
\textit{b}, et \textit{C}, \textit{c}, et celeritas \textit{B} sit \textit{a}, ipsius 
\edtext{\textit{C} \textit{e}. erit}{%
\lemma{\textit{C} \textit{e}.}%
\Bfootnote{%
\textit{(1)}~sit %
\textit{(2)}~erit~\textit{L}}}
potentia ante concursum\indexS{concursus}\indexS{potentia ante concursum} $ba^2+ce^2$. celeritas post concursum\indexS{concursus}\indexS{celeritas post concursum} erit
\edtext{$\overline{ab+ec}:\overline{b+c}$, cujus quadratum}{%
\lemma{$\overline{ab+ec}:\overline{b+c}$}%
\Bfootnote{%
\textit{(1)}~cujus celeritas %
\textit{(2)}~, itaque %
\textit{(3)}~cujus quadratum~\textit{L}}}
ductum\indexS{quadratum celeritatis ductum in corpus} in
\edtext{\lbrack$b+c$\rbrack}{%
\lemma{}%
\Bfootnote{%
$a+b$ %
\textit{L ändert Hrsg.}}}
dabit: $a^2b^2+e^2c^2+2abce,:\overline{b+c}$ quae erit potentia\indexS{potentia post ictum}
\edtext{post ictum,\indexS{ictus} ut autem quaeratur differentia a}{%
\lemma{post ictum,}%
\Bfootnote{%
\textit{(1)}~ea autem multo minor est quam %
\textit{(2)}~ut autem quaeratur differentia a~\textit{L}}}
$ba^2+ce^2$ potentia\indexS{potentia ante ictum}
\edtext{ante ictum,\indexS{ictus} sit enim ipsius $a^2b^2+e^2c^2+2abce$ differentia a $b^2a^2}{%
\lemma{ante ictum,}%
\Bfootnote{%
\textit{(1)}~et differentia est %
\textit{(2)}~sit enim \bars{ipsius \textit{erg.}} $a^2b^2+e^2c^2+2abce$ differentia a %
\textit{(a)}~$ba^3+ace$ %
\textit{(b)}~$b^2a^2$~\textit{L}}}
+bce^2+bca^2+c^2e^2$, dividenda per $b+c$,  
\edtext{quae differentia divisa erit}{%
\lemma{quae differentia}%
\Bfootnote{%
\textit{(1)}~erit %
\textit{(2)}~divisa erit~\textit{L}}}
$e^2+a^2-2ae,\cdot\,bc:\overline{b+c}$. Unde patet quia $e^2+a^2$ est semper majus quam $2ae$, 
\edtext{etiam potentiam\indexS{potentia post ictum}}{%
\lemma{etiam}%
\Bfootnote{%
\textit{(1)}~vim %
\textit{(2)}~potentiam~\textit{L}}}
post ictum\indexS{ictus} 
\edtext{semper esse}{%
\lemma{semper}%
\Bfootnote{%
\textit{(1)}~erit %
\textit{(2)}~esse~\textit{L}}}
minorem quam ante ictum\indexS{ictus}\indexS{potentia ante ictum} et 
\edtext{calculo inveniri}{%
\lemma{}%
\Bfootnote{%
calculo \bars{prioribus simili \textit{gestr.}} inveniri~\textit{L}}}
poterit, quando differentia sit 
\edtext{maxima. Et jam tum patet quo}{%
\lemma{maxima.}%
\Bfootnote{%
\textit{(1)}~Sane quia \textit{e} est minor, hoc %
\textit{(2)}~Et jam tum patet quo~\textit{L}}}
celeritas ipsius \textit{e} minor assumitur, et hoc $e^2+a^2-2ae$ esse 
\edtext{majorem. Quod autem ad $bc:\overline{b+c}$ attinet sit}{%
\lemma{majorem}%
\Bfootnote{%
\textit{(1)}~, et rursus, patet quo %
\textit{(2)}~. Quod autem ad $bc:\overline{b+c}$ attinet %
\textit{(a)}~seu ad $b:c$ %
\textit{(b)}~sit~\textit{L}}}
$b=1$ et fiet $bc:\overline{b+c}=c:\overline{1+c}$, quod quaeritur omnium possibilium
\edtext{maximum,\indexS{maximum} seu}{%
\lemma{maximum,}%
\Bfootnote{%
\textit{(1)}~sit \textit{c} %
\textit{(2)}~seu~\textit{L}}}
$\displaystyle\frac{1+c}{c}$ omnium possibilium 
\edtext{minimum\indexS{minimum} $=z$. Assumsimus}{%
\lemma{minimum $=z$.}%
\Bfootnote{%
\textit{(1)}~Patet autem %
\textit{(2)}~Assumsimus~\textit{L}}}
autem \textit{c} esse majus quam \textit{b}. 
Itaque si \textit{c} sit infinitum, \textit{z} erit 1, si aequale ipsi \textit{b} erit $z=2$, 
\edtext{sed calculo differentiali\indexS{calculus differentialis}}{%
\lemma{sed}%
\Bfootnote{%
\textit{(1)}~quorum %
\textit{(2)}~calculo %
\textit{(a)}~Maximor %
\textit{(b)}~differentiali~\textit{L}}}
fiet: 
\edtext{$\lbrack b\rbrack=c+1$}{%
\lemma{}%
\Bfootnote{%
$c=c+1$ %
\textit{L ändert Hrsg.}}}
seu $c=0$.
\edtext{Unde cum}{%
\lemma{Unde}%
\Bfootnote{%
\textit{(1)}~patet absolute aeq. %
\textit{(2)}~cum~\textit{L}}}
sit una tantum radix ipsius, \textit{c}, 
\edtext{patet maximum\indexS{maximum}}{%
\lemma{patet}%
\Bfootnote{%
\textit{(1)}~Mini %
\textit{(2)}~mi %
\textit{(3)}~in %
\textit{(4)}~Minimum %
\textit{(5)}~maximum~\textit{L}}}
haberi cum $c=0$, minimum\indexS{minimum} autem nullum esse, sed \textit{c} posse augeri in infinitum, seu maximum\indexS{maximum} esse cum \textit{c} est infinitum limites autem sumi ex datis, 
\edtext{quod debet esse et $c \groesser b$.}{%
\lemma{quod}%
\Bfootnote{%
\textit{(1)}~\textit{b} %
\textit{(2)}~$a:e$ %
\textit{(3)}~debet esse %
\textit{(a)}~major %
\textit{(b)}~et $c \groesser b$.~\textit{L}}}
et $a:e \groesser b:c$. 
\edtext{Sumatur ergo celeritas corporis \textit{B} utcunque magna respectu ejus quam habet \textit{C} et magnitudo corporis \textit{C} utcunque magna respectu corporis \textit{B}, unde semper maneret ratio celeritatum\indexS{ratio celeritatum} major quam magnitudinum.\indexS{ratio magnitudinum}}{%
\lemma{}%
\Bfootnote{%
Sumatur ergo %
\textit{(1)}~utcunque %
\textit{(2)}~corpus \textit{C} %
\textit{(3)}~celeritas corporis \textit{B} utcunque magna respectu %
\textit{(a)}~corporis %
\textit{(b)}~ejus \lbrack...\rbrack\ ratio %
\textit{(aa)}~magnit %
\textit{(bb)}~celeritatum major quam magnitudinum. \textit{erg.~L}}} 
\pend
\pstart
\edtext{Caeterum in}{%
\lemma{}%
\Bfootnote{%
Caeterum
\bars{\textit{(1)}~hoc \textit{(2)}~haec \textit{streicht Hrsg.}} in~\textit{L}}}
universum in regulis 
\edtext{Cartesii\indexN{\textls{Descartes} (Cartesius, des Cartes), René 1596\textendash1650}\indexS{regulae motus Cartesii}
haec observo: primum non esse}{%
\lemma{Cartesii}%
\Bfootnote{%
\textit{(1)}~observo. Non esse %
\textit{(2)}~haec observo: primum non esse~\textit{L}}}
completas\lbrack,\rbrack\ 
\edlabel{35_10_12_038_8a}%
omissi enim casus valde notabiles, v.\,g. cum corpus unum attingit aliud antecedens 
celeritate magnitudinibus reciproca,\indexS{celeritas magnitudini reciproca} sed in universum 
quando duo corpora inaequalia\indexS{corpora inaequalia} inaequalibus celeritatibus\indexS{celeritates inaequales} sibi obviam eunt.\edlabel{35_10_12_038_8b}	%
Deinde 
\edtext{peccat graviter}{%
\lemma{peccat}%
\Bfootnote{%
\textit{(1)}~gravissime %
\textit{(2)}~graviter~\textit{L}}}
quod non considerat, ictum\indexS{ictus} duorum corporum esse proportionalem celeritatibus 
\edtext{e\lbrack o\rbrack rum}{%
\lemma{}%
\Bfootnote{%
earum %
\textit{L ändert Hrsg.}}}
non absolutis,\indexS{celeritates corporum absolutae} sed 
\edtext{respectivis,\indexS{celeritates corporum absolutae} et si eadem}{%
\lemma{respectivis,}%
\Bfootnote{%
\textit{(1)}~eundemque %
\textit{(2)}~et si eadem~\textit{L}}}
celeritate sibi appropinquant, eundem esse, in quacunque demum parte sit motus absolutus\indexS{motus absolutus} sive verus.\indexS{motus verus} 
\edtext{Sed gravissima omnium peccata\indexS{peccata gravissima Cartesii} sunt duo, quod}{%
\lemma{Sed}%
\Bfootnote{%
\textit{(1)}~gravissimum omnium peccatum est %
\textit{(2)}~gravissima omnium peccata sunt duo, 
\bars{primum \textit{gestr.}} quod~\textit{L}}}
effectus est inaequalis causae,\indexS{effectus inaequalis causae} 
ut jam in aliquot exemplis ostendimus, et quod mutatione in datis 
seu suppositis\indexS{mutatio in datis seu suppositis} 
facta utcunque exigua eventus\indexS{eventus} tamen in immensum discrepant,\indexS{discrepantia in eventibus} 
et levissima mutatione quae sit data quavis minor\indexS{mutatio data quavis minor} transitur 
in oppositam, quod naturae oppositum est, et ostendit regulas\indexS{regulae motus Cartesii} esse male concinnatas. 
Et incompatibile est cum calculo.
Si enim res ad calculum 
\edtext{generalem redacta}{%
\lemma{generalem}%
\Bfootnote{%
\textit{(1)}~redigi possit, necesse est casus %
\textit{(2)}~redacta~\textit{L}}}
sit, semper casus quietis debet posse comprehendi in casu motus, supponendo motum infinite parvum.
\pend
